% GENERATED FILE. Edit canonical lessons, then run npm run generate:books.
% canonical-source-hash: fnv1a64:3bd6af51c88cb81f
% canonical-lessons: GE-C519-galerie, GE-C519-ausstellen, GE-C519-skulptur, GE-C519-bildhauer, GE-C519-kunstwerk

\chapter{In der Galerie}
\label{ch:ge-in-der-galerie}
\hlchaptermodality{519}

\begin{chapteropening}
\textbf{By the end of this chapter:} \emph{I can talk about art exhibitions.}

Complete the last lesson of chapter 519: I can talk about art exhibitions.
\end{chapteropening}

\section[die Galerie]{die Galerie — gallery}

\label{lesson:GE-C519-galerie}



\begin{quote}
Before the new one: say the German for easel, then the German for watercolour.
\end{quote}

\subsection*{You'll want to know: die Galerie}
\textbf{die Galerie} — \textquotedblleft{}gallery\textquotedblright{}.

\begin{quote}
In der Galerie hängen Bilder von jungen Künstlern.
\end{quote}

\begin{quote}
\emph{Pictures by young artists hang in the gallery.}
\end{quote}

Plural: \emph{die Galerien}.

\begin{cousinweb}
\emph{Galerie} came into German from Italian \emph{galleria}, through French \emph{galerie}; the older origin of the Italian word is not settled, and EtymDB lists the German word without a source.
\end{cousinweb}

\begin{grammarlens}[title={What you've built}]
One more word for telling, explaining and describing at B1.
\end{grammarlens}

\subsection*{Your turn}
\begin{itemize}
\raggedright
  \item \emph{Say it:} \emph{die Galerie}
  \item \emph{Say it:} \emph{die Galerie}, once more
  \item \emph{Say it:} a sentence of your own with \emph{die Galerie}
  \item \emph{Recall:} say the German for easel, then the German for watercolour, then say \emph{die Galerie} again
\end{itemize}

\begin{tcolorbox}[breakable,colback=teal!4,colframe=teal!35!black,title={Before you move on}]
What does die Galerie mean? (\textquotedblleft{}Gallery\textquotedblright{}.) Say it once more.
\end{tcolorbox}
\section[ausstellen]{ausstellen — to exhibit; to issue}

\label{lesson:GE-C519-ausstellen}



\begin{quote}
Before the new one: say the German for watercolour, then the German for gallery.
\end{quote}

\subsection*{You'll want to know: ausstellen}
\textbf{ausstellen} — \textquotedblleft{}to exhibit; to issue\textquotedblright{}.

\begin{quote}
Das Museum stellt Fotos aus den Sechzigerjahren aus.
\end{quote}

\begin{quote}
\emph{The museum is exhibiting photos from the sixties.}
\end{quote}

Separable: \emph{stellt aus}, \emph{hat ausgestellt}. Also to issue a document: \emph{einen Pass ausstellen}.

\begin{cousinweb}
\emph{ausstellen} is built on \emph{aus} \textquotedblleft{}out\textquotedblright{} and \emph{stellen} \textquotedblleft{}to set, to place\textquotedblright{} (EtymDB), inherited from Old High German \emph{stellen}, from Proto-Germanic \emph{stallijaną}.
\end{cousinweb}

\begin{grammarlens}[title={What you've built}]
One more word for telling, explaining and describing at B1.
\end{grammarlens}

\subsection*{Your turn}
\begin{itemize}
\raggedright
  \item \emph{Say it:} \emph{ausstellen}
  \item \emph{Say it:} \emph{ausstellen}, once more
  \item \emph{Say it:} a sentence of your own with \emph{ausstellen}
  \item \emph{Recall:} say the German for watercolour, then the German for gallery, then say \emph{ausstellen} again
\end{itemize}

\begin{tcolorbox}[breakable,colback=teal!4,colframe=teal!35!black,title={Before you move on}]
What does ausstellen mean? (\textquotedblleft{}To exhibit; to issue\textquotedblright{}.) Say it once more.
\end{tcolorbox}
\section[die Skulptur]{die Skulptur — sculpture}

\label{lesson:GE-C519-skulptur}



\begin{quote}
Before the new one: say the German for gallery, then the German for to exhibit.
\end{quote}

\subsection*{You'll want to know: die Skulptur}
\textbf{die Skulptur} — \textquotedblleft{}sculpture\textquotedblright{}.

\begin{quote}
Vor dem Rathaus steht eine moderne Skulptur aus Metall.
\end{quote}

\begin{quote}
\emph{In front of the town hall there's a modern metal sculpture.}
\end{quote}

Plural: \emph{die Skulpturen}.

\begin{cousinweb}
\emph{Skulptur} is a learned borrowing from Latin \emph{sculptūra}, from \emph{sculpere} \textquotedblleft{}to carve\textquotedblright{} (Lewis \& Short).
\end{cousinweb}

\begin{grammarlens}[title={What you've built}]
One more word for telling, explaining and describing at B1.
\end{grammarlens}

\subsection*{Your turn}
\begin{itemize}
\raggedright
  \item \emph{Say it:} \emph{die Skulptur}
  \item \emph{Say it:} \emph{die Skulptur}, once more
  \item \emph{Say it:} a sentence of your own with \emph{die Skulptur}
  \item \emph{Recall:} say the German for gallery, then the German for to exhibit, then say \emph{die Skulptur} again
\end{itemize}

\begin{tcolorbox}[breakable,colback=teal!4,colframe=teal!35!black,title={Before you move on}]
What does die Skulptur mean? (\textquotedblleft{}Sculpture\textquotedblright{}.) Say it once more.
\end{tcolorbox}
\section[der Bildhauer]{der Bildhauer — sculptor}

\label{lesson:GE-C519-bildhauer}



\begin{quote}
Before the new one: say the German for to exhibit, then the German for sculpture.
\end{quote}

\subsection*{You'll want to know: der Bildhauer}
\textbf{der Bildhauer} — \textquotedblleft{}sculptor\textquotedblright{}.

\begin{quote}
Der Bildhauer arbeitet mit Holz und Stein.
\end{quote}

\begin{quote}
\emph{The sculptor works with wood and stone.}
\end{quote}

Plural: \emph{die Bildhauer}. A woman is \emph{die Bildhauerin}.

\begin{cousinweb}
\emph{Bildhauer} is built on \emph{Bild} \textquotedblleft{}image\textquotedblright{}, from Proto-Germanic \emph{biliþją}, and \emph{hauen} \textquotedblleft{}to hew\textquotedblright{}, inherited from Old High German \emph{houwan} (EtymDB): someone who hews images.
\end{cousinweb}

\begin{grammarlens}[title={What you've built}]
One more word for telling, explaining and describing at B1.
\end{grammarlens}

\subsection*{Your turn}
\begin{itemize}
\raggedright
  \item \emph{Say it:} \emph{der Bildhauer}
  \item \emph{Say it:} \emph{der Bildhauer}, once more
  \item \emph{Say it:} a sentence of your own with \emph{der Bildhauer}
  \item \emph{Recall:} say the German for to exhibit, then the German for sculpture, then say \emph{der Bildhauer} again
\end{itemize}

\begin{tcolorbox}[breakable,colback=teal!4,colframe=teal!35!black,title={Before you move on}]
What does der Bildhauer mean? (\textquotedblleft{}Sculptor\textquotedblright{}.) Say it once more.
\end{tcolorbox}
\section[das Kunstwerk]{das Kunstwerk — work of art}

\label{lesson:GE-C519-kunstwerk}



\begin{quote}
Before the new one: say the German for sculpture, then the German for sculptor.
\end{quote}

\subsection*{You'll want to know: das Kunstwerk}
\textbf{das Kunstwerk} — \textquotedblleft{}work of art\textquotedblright{}.

\begin{quote}
Für mich ist dieser Garten ein echtes Kunstwerk.
\end{quote}

\begin{quote}
\emph{To me this garden is a real work of art.}
\end{quote}

Plural: \emph{die Kunstwerke}.

\begin{cousinweb}
\emph{Kunstwerk} is built on \emph{Kunst} \textquotedblleft{}art, skill\textquotedblright{} and \emph{Werk} \textquotedblleft{}work\textquotedblright{}, inherited from Old High German \emph{werc}, from Proto-Germanic \emph{werką} (EtymDB).
\end{cousinweb}

\begin{grammarlens}[title={What you've built}]
One more word for telling, explaining and describing at B1.
\end{grammarlens}

\subsection*{Your turn}
\begin{itemize}
\raggedright
  \item \emph{Say it:} \emph{das Kunstwerk}
  \item \emph{Say it:} \emph{das Kunstwerk}, once more
  \item \emph{Say it:} a sentence of your own with \emph{das Kunstwerk}
  \item \emph{Recall:} say the German for sculpture, then the German for sculptor, then say \emph{das Kunstwerk} again
\end{itemize}

\begin{tcolorbox}[breakable,colback=teal!4,colframe=teal!35!black,title={Before you move on}]
What does das Kunstwerk mean? (\textquotedblleft{}Work of art\textquotedblright{}.) Say it once more.
\end{tcolorbox}
